\section*{章末問題}
\begin{enumerate}[label=\textbf{\thechapter.\arabic*},leftmargin=*]
 \item\label{問:定積分リーマン和} 区間 $[0,2]$ を $n$ 等分し,各小区間の右端を使ったリーマン和から $\int_0^2x^2\dd x$ を求めよ.幅と代表点を明記せよ.
 \item\label{問:定積分性質} 連続関数 $f,g$ が
 \[
  \int_0^1f(x)\dd x=2,\qquad
  \int_1^3f(x)\dd x=-1,\qquad
  \int_0^3g(x)\dd x=4
 \]
 を満たす.\,$\int_0^3f(x)\dd x$ と $\int_3^0\{2f(x)-3g(x)\}\dd x$ を求めよ.
 \item\label{問:定積分と面積} $\int_{-1}^2(x-1)\dd x$ と,直線 $y=x-1$,\,$x$ 軸,直線 $x=-1,x=2$ で囲まれる部分の面積を求めよ.符号が変わる点で区間を分けて説明せよ.
 \item\label{問:基本定理計算} 基本定理を使って次を求めよ.
 \[
  \text{(1)}\quad\int_{-1}^2(x^2+2x)\dd x,\qquad
  \text{(2)}\quad\int_0^{\frac\pi2}(\sin x+\cos x)\dd x,
 \]
 \[
  \text{(3)}\quad\int_1^{\e}\left(\frac2x+x\right)\dd x.
 \]
 \item\label{問:面積関数増減} $A(x)=\int_1^x(t-2)(t+1)\dd t$ とする.\,$A'(x)$ を求め,区間 $[-2,3]$ における $A$ の最大値・最小値と,それらをとる $x$ を求めよ.
 \item\label{問:変数上下端} $G(x)=\int_{\cos x}^{\sin x}\sqrt{1+t^2}\dd t$ を微分せよ.下端から生じる項の符号に注意せよ.
 \item\label{問:積分の評価} $0\le x\le1$ で $1\le\sqrt{1+x^2}\le1+\frac{x^2}{2}$ が成り立つことを示し,それを用いて
 \[
  1\le\int_0^1\sqrt{1+x^2}\dd x\le\frac76
 \]
 を証明せよ.平方根の原始関数を求める必要はない.
 \item\label{問:積分方程式基本} 次の条件を満たす連続関数を求めよ.存在しない場合は,その理由を示せ.
 \begin{enumerate}[label=(\arabic*),leftmargin=*]
  \item すべての実数 $x$ で $f(x)=3-\int_0^x f(t)\dd t$.
  \item すべての $x\in[0,1]$ で $g(x)=x+\int_0^1g(t)\dd t$.
 \end{enumerate}
\end{enumerate}

\smallskip\noindent 解答は\pageref{解答:第4章}ページ.

